\section{Finite accounting of chiral Weyl energy}

We work on one common source sequence. The base and all bubbles contribute
in one chirality. All measures in the accounting theorems have finite
mass. For the geometric application the source measure is
$\mu_i=|W^+_{g_i}|_{\mathrm{HS}}^2\,dV_{g_i}$. A smooth compactified
base can carry positive energy after all bubble energies vanish.

\begin{lemma}[Finite accounting with a strict residual]
\label{lem:finite-bubble-accounting}
\lean{SectionsAndConnections.noPositiveBubbleOfFiniteAccounting}
\leanok
Put $p=4\pi^2$. Let $J$ be finite, $\delta>0$ and $p\le e\le2p-\delta$.
Suppose $e=b+\sum_{j\in J}x_j$, where $b,x_j\ge0$. Each of $b$ and
the $x_j$ is zero or strictly larger than $p$. If $e-x_j>0$ whenever
$x_j>0$, then every $x_j=0$ and $e=b$.
\end{lemma}
\begin{proof}
\leanok
For a positive $x_j$, the residual $e-x_j$ is less than $p$. It is
the sum of $b$ and all the other bubble energies. Each is nonnegative
and strictly below $p$, so each is zero. This contradicts $e-x_j>0$.
\end{proof}

\begin{lemma}[Source and component budgets in scalar form]
\label{lem:scalar-component-budgets}
\lean{FirstOrderTwistor.sourceBandScalarBudget, FirstOrderTwistor.sourceCurvatureBudget, FirstOrderTwistor.strictComponentResidual, FirstOrderTwistor.uniformComponentResidual}
\leanok
Write $Y_{\rm rd}=8\sqrt6\,\pi$. Given the source Euler identity
$Y^2=48(8\pi^2-E)+12D$, $0\le Y\le Y_{\rm rd}$, $D\ge0$,
and $E\le8\pi^2-\delta$ with $\delta>0$, one has $Y\ge\sqrt{48\delta}$
and $D\le4E$. If the full curvature integral is $8E+2D+Y^2/6$, it is
at most $192\pi^2-16\delta$. For a component with energy $x$, suppose
the limiting Euler identity holds at $(e,y,D)$ and $4x=d\le D$. Then
\[
 e-x\ge (Y_{\rm rd}^2-y^2)/48.
\]
If $0\le y\le\widehat Y<Y_{\rm rd}$ the right hand side is positive.
If $\widehat Y\le Y_{\rm rd}-d_\delta$ for a supplied
$0<d_\delta<Y_{\rm rd}$, the lower bound is at least
$[Y_{\rm rd}^2-(Y_{\rm rd}-d_\delta)^2]/48$.
\end{lemma}
\begin{proof}
\leanok
Substitute the Euler identity and $Y_{\rm rd}^2=384\pi^2$. On the
nonnegative real axis, squaring preserves the Yamabe comparisons.
\end{proof}

\begin{lemma}[Counting ends from the numerical end budget]
\label{lem:scalar-ale-end-count}
\lean{FirstOrderTwistor.oneALEEndOfSubcriticalBounds, FirstOrderTwistor.q8EtaContradictionOfSubcriticalBand}
\leanok
In the notation of \S11.3 let $k\ge1$ count the end orders $q_j\ge1$,
all bounded by $Q$. Suppose $e\le8\pi^2-\delta$, $\delta>0$,
$e_+,e_-\ge0$, and
\[
 (8\pi^2-e)+(e_++e_-)/2+4\pi^2(k-2+\textstyle\sum_j1/q_j)
 \le 8\pi^2/Q.
\]
Then $k=1$. If both chiral energies belong to $[0,8\pi^2-\delta]$,
the $Q_8$ signature defect $|e_+-e_-|=9\pi^2$ is impossible.
\end{lemma}
\begin{proof}
\leanok
If $k\ge2$, insert $\sum_j1/q_j\ge k/Q$ and subtract $8\pi^2/Q$.
The difference is at least $\delta+4\pi^2(k-2)(1+1/Q)>0$.
The two nonnegative chiral energies differ by less than $9\pi^2$.
\end{proof}

\begin{lemma}[Energy on a fixed comparison domain]
\label{lem:fixed-domain-density}
\lean{FirstOrderTwistor.energyMeasure_real_range_of_embedding, FirstOrderTwistor.tendstoImageEnergyOfUniformPullback, FirstOrderTwistor.chiralHSDensity_integrableOn_and_tendsto}
\leanok
For a measurable embedding $\phi_i$ of a fixed finite-volume domain,
suppose the pullback volume has nonnegative density $J_i$ relative to
$d\nu$. For an integrable nonnegative source density $f_i$ on its image,
the image mass is $\int J_i(f_i\circ\phi_i)\,d\nu$.
Uniform convergence of these densities to an integrable nonnegative
density gives convergence of image masses. For given chiral matrices
$W_i$ in oriented orthonormal frames, almost-everywhere convergence of
the nine entries and $J_i$ also gives convergence of
$\int J_i\sum_{a,b}(W_i)_{ab}^2\,d\nu$ under measurability and a common
integrable majorant.
\end{lemma}
\begin{proof}
\leanok
Transport measure by the embedding and use the formula for integration
against a measure with density. Apply the finite-volume uniform bound
or dominated convergence to the finite sum of squares.
\end{proof}

\begin{lemma}[The mass of the complement]
\label{lem:source-remainder-split}
\lean{FirstOrderTwistor.measureReal_univ_eq_sum_add_compl, FirstOrderTwistor.measureReal_le_sum_add_uncovered}
\leanok
For finitely many measurable $C_j$ with pairwise null intersections,
put $R=(\bigcup C_j)^c$. Then $\mu(X)=\sum_j\mu(C_j)+\mu(R)$.
If $\mu(S_k)\le d_k$ for finitely many strips, then
\[
 \mu(R)\le\sum_k d_k+\mu(R\setminus\textstyle\bigcup_k S_k).
\]
\end{lemma}
\begin{proof}
\leanok
Use finite additivity up to null sets for the first equation and
subadditivity on the covered portion for the second.
\end{proof}

\begin{theorem}[The energy defect and equality under a strip estimate]
\label{thm:covered-measure-accounting}
\uses{lem:source-remainder-split}
\lean{FirstOrderTwistor.TruncationFamily.sum_componentEnergy_le, FirstOrderTwistor.TruncationFamily.tendsto_limsup_remainder, FirstOrderTwistor.energyIdentityOfCoveredRemainder}
\leanok
Suppose $\mu_i(X)\to e$. Each of finitely many component measures
$\nu_j$ has increasing truncations $K_L^j$ exhausting its space up to
a null set. For fixed $L$, the source images $C_{Li}^j$ are eventually
measurable and pairwise disjoint up to null sets, with
$\mu_i(C_{Li}^j)\to\nu_j(K_L^j)$. With $R_{Li}=(\bigcup_j C_{Li}^j)^c$
and $s=\sum_j\nu_j(X_j)$, one has
\[
 s\le e,\qquad \lim_{L\to\infty}\limsup_{i\to\infty}\mu_i(R_{Li})=e-s .
\]
Suppose also that each strip mass is at most its estimated cost, and
that both the sum of costs and the uncovered mass tend to zero in the
following sense. For every $\varepsilon>0$ and all sufficiently large
$L$, each quantity is eventually at most $\varepsilon$ as $i\to\infty$.
Then $e=s$. The finite strip family may depend on $L$ and $i$.
\end{theorem}
\begin{proof}
\uses{lem:source-remainder-split}
\leanok
At fixed $L$, take the limit of the exact mass split. The finite sum of
truncated masses then converges to $s$ by continuity of measure from
below. Bounding each of the two errors by $\varepsilon/2$ gives $e=s$.
\end{proof}

\begin{corollary}[Bubble energies with supplied geometric budgets]
\label{cor:covered-bubble-energies}
\uses{thm:covered-measure-accounting,lem:scalar-component-budgets,lem:finite-bubble-accounting}
\lean{FirstOrderTwistor.noPositiveBubbleOfCoveredRemainder, FirstOrderTwistor.bubbleEnergyZeroOfRicciYamabeBounds}
\leanok
With one distinguished base and the band and gap hypotheses of
Lemma~\ref{lem:finite-bubble-accounting}, the covered-remainder hypotheses
give zero mass for every bubble and $e=e_0$ if each positive bubble has
strict residual. For the same conclusion it suffices to supply the
limiting Euler identity, $4e_j\le D$, and
$0\le y\le\widehat Y_j<Y_{\rm rd}$ for each positive bubble.
\end{corollary}
\begin{proof}
\uses{thm:covered-measure-accounting,lem:scalar-component-budgets,lem:finite-bubble-accounting}
\leanok
Derive equality by exhaustion. The component budget supplies the strict
residual in the second formulation, so the finite accounting lemma applies.
\end{proof}

\begin{lemma}[Sequence limits of the supplied annular bounds]
\label{lem:annular-bound-limits}
\lean{FirstOrderTwistor.noNeckOfSuppliedCharts, FirstOrderTwistor.noNeckOfSuppliedField, FirstOrderTwistor.annularSmallness_iff_tendsto_limsup, FirstOrderTwistor.smallSumOfFiniteStripErrors}
\leanok
Given the numerical estimate of Theorem~\ref{thm:one-chart-energy}, a
uniform bound on $K,\mathfrak B$ and iterated smallness of $\varepsilon$
give iterated smallness of the annular energy. The same holds for
the intrinsic bound $8\Lambda\mathfrak B+4\|W\|_2\mathfrak d$
when $\Lambda,\|W\|_2$ are uniformly bounded and both error inputs tend
to zero in the iterated sense. For nonnegative energies eventually
bounded above at each fixed large $L$, this is equivalent to the
iterated limsup being zero. Individual estimates sum over a fixed
finite strip family.
\end{lemma}
\begin{proof}
\leanok
Bound the coefficients independently of large $L$ and sufficiently
large $i$ at that $L$. Apply the defining eventual bounds for the
inner limsup and sum over finitely many strips.
\end{proof}

\begin{lemma}[Covered portion of a single end]
\label{lem:diagonal-covered-end}
\lean{FirstOrderTwistor.exists_slow_diagonal, FirstOrderTwistor.exists_slow_diagonal_with_vanishing_annularMeasure}
\leanok
Suppose each fixed truncation eventually passes its comparison test.
Suppose its annular mass between $A$ and $m$ is at most the mass of
a model annulus in the tail beyond $A$ plus an error tending to zero
with $m$. If the model tail mass tends to zero with $A$, there is
a choice $m(i)\to\infty$ passing every test through $m(i)$ eventually
and giving iterated smallness of this covered annular mass.
\end{lemma}
\begin{proof}
\leanok
Choose the largest $m\le i$ for which the tests through $m$ pass.
For each fixed $m$ this prefix eventually passes. Compare at the
selected endpoint and then exhaust the inner tail. Energy up to a
different prescribed outer scale requires a separate estimate.
\end{proof}
